\section{Limitazioni}
\label{sec:limitazioni}

\subsection{Scalabilità durante l'esecuzione}

I Cloud Worker e il Global Aggregator mantengono in memoria lo stato delle
finestre in elaborazione. Tale stato non viene migrato tra Worker e non è
previsto un meccanismo di checkpoint che consenta di trasferirlo durante una
run.

Per questo motivo la scalabilità viene valutata configurando il numero di
Cloud Worker prima dell'avvio del replay e confrontando esecuzioni distinte.
Un rebalance del consumer group dopo l'inizio dell'elaborazione invalida la
run e i Worker non riprendono una precedente esecuzione a partire da offset
già committati. Il sistema permette quindi di studiare l'aumento del
parallelismo tra esecuzioni differenti, ma non implementa una variazione
dinamica del numero di Worker durante la stessa run.

\subsection{Completezza e rappresentatività}

La telemetria Simulator--Edge utilizza MQTT con QoS~0 e attraversa code di
capacità limitata. Per evitare che eventuali perdite vengano interpretate
come un miglioramento prestazionale, nella campagna di scalabilità vengono
considerate valide soltanto le esecuzioni che mantengono il 100\% degli
eventi, senza drop o contributi late.

Negli esperimenti bounded, invece, i contributi late costituiscono un effetto
previsto della politica event-time studiata. Una sorgente che rimane indietro
oltre lo skew configurato può essere superata dalla frontiera; i contributi
che raggiungono successivamente una finestra già finalizzata vengono
registrati come late e scartati. Eventuali drop nei Simulator, negli Edge o
nell'emulazione di rete rimangono invece condizioni di invalidità della run.

Il workload utilizza un solo tipo di dispositivo, i sensori BME280, e
considera le osservazioni relative a un solo mese. I 150 sensori selezionati
sono distribuiti in modo quasi uniforme tra i 13 Edge Site, ma il numero di
eventi prodotto dalle diverse zone non è necessariamente uniforme. Le medie
sono inoltre calcolate a partire dal numero di campioni validi, per cui
contributi contenenti un numero maggiore di osservazioni hanno un peso
maggiore nel risultato globale.

I Cloud Worker condividono il medesimo host Workers e sono eseguiti con i
limiti di CPU e memoria definiti dal profilo sperimentale. I risultati
descrivono pertanto il comportamento del deployment e del workload analizzati
e non vengono generalizzati automaticamente ad altri dataset, algoritmi più
costosi o configurazioni infrastrutturali differenti.

\section{Conclusioni}
\label{sec:conclusioni}

La campagna AWS mostra che l'aumento del parallelismo del livello Cloud riduce
la latenza di emissione mantenendo invariati workload e risultati. Passando da
uno a quattro Cloud Worker, la mediana della \emph{global emission latency}
diminuisce di circa il 40\% e il p95 di circa il 32\%. Con sei Worker viene
raggiunto il massimo parallelismo consentito dalle sei partizioni Kafka, ma
non si osserva un ulteriore miglioramento: la mediana rimane sostanzialmente
invariata, mentre p95 e p99 risultano leggermente superiori rispetto a W4.
I risultati mostrano quindi un beneficio del parallelismo fino a tre-quattro
Worker e rendimenti decrescenti oltre tale configurazione per il workload e
il deployment analizzati.

La campagna Edge--Cloud evidenzia invece il compromesso tra latenza e
completezza introdotto dalla politica bounded. Con 50 ms di ritardo aggiunto,
la configurazione conservative mantiene il 100\% degli eventi con un p95 pari
a 3.87 s, mentre la configurazione bounded riduce il p95 a 0.83 s con una
completezza dell'89.12\%. Con il collegamento di \code{edge-8} limitato a
140 kbit/s, la configurazione conservative mantiene il 100\% degli eventi con
un p95 pari a 37.82 s, mentre la configurazione bounded riduce il p95 a
1.58 s con una completezza dell'82.44\%.

Nelle configurazioni bounded la riduzione della completezza è spiegata dai
contributi registrati come late, mentre nelle esecuzioni considerate non sono
presenti drop nei Simulator, negli Edge o nell'emulazione di rete. Il sistema
permette quindi di osservare e quantificare il compromesso tra il tempo di
attesa di una sorgente rallentata e la completezza dei risultati prodotti.